\section{Conclusions}
\label{sec:conclusions}

We addressed validator-set sizing for permissioned BFT when no large testbed is
available. The identifiability theorem shows that a single-factor exponent
converges to $\gamma\lambda-\beta$, so a positive fit is consistent with a load
generator scaled alongside the network and explains the discrepancy between our
earlier predictions~\cite{peniak2026a,peniak2026b}. The \RNM{} protocol makes
the required factorial design usable on one commodity host; where exact
$\q$-invariance fails, the law is reported at the reference quota. Coverage-
validated prediction intervals replace unchecked point accuracy, and coupling
the throughput law with a Hoeffding-type detection bound yields a
chance-constrained program whose feasible set is the closed interval
$[\nmin,\nmax]$, with absolute demand entering through $\kappaScale$. Of the
four results, the first --- a statement about estimation --- is the most likely
to outlive the specific measurements. The full campaign runs in a public CI
workflow on one four-core runner.

\paragraph{Future work.}
Equivocation-capable instrumentation; multi-client validation; latency/finality
constraints in the window; online recalibration of~$\beta$.

\subsection*{Conflict of interest}
The authors declare no conflict of interest.

\subsection*{Funding}
The research was conducted without financial support.

\subsection*{Data availability}
Workflow definition, generators, fault-injection and analysis code, and raw
artefacts are available at
\url{https://github.com/bpenyak/paper_3}.

\subsection*{Use of artificial intelligence}
AI-assisted tools were used for language editing and typesetting. All
scientific results, proofs, design and analysis were produced by the authors.

\subsection*{Contributions of authors}
\textbf{Peniak B.\,O.}: conceptualization, methodology, formal analysis,
software, investigation, data curation, writing --- original draft,
visualization.
\textbf{Liubinskyj B.\,B.}: supervision, validation, writing --- review and
editing.
